\section{2차 전기 시스템}

2장에서는 긴 선로가 왜 임피던스라는 언어를 요구했는지 보았고, 3장에서는 그 언어를 LTI 시스템과 전달함수로 읽는 법을 정리했다.
긴 선로는 위치마다 조금씩 다른 저항, 인덕턴스, 커패시턴스, 누설 성분을 가진 복잡한 시스템이다.
여기서는 그 전체 선로를 바로 풀지 않고, 더 작은 2차 예제인 RLC 회로부터 살펴본다.
작은 RLC 회로를 이해하면, 긴 선로에서 왜 저항 하나로는 부족했는지 훨씬 쉽게 보인다.

\subsection{진동이란 무엇인가?}

진동(oscillation)은 어떤 물리량이 평형점을 중심으로 반복해서 커졌다 작아지는 운동이다.
추를 잡아당겼다가 놓으면 좌우로 흔들리고, 기타 줄을 튕기면 위아래로 떨리며, RLC 회로에서는 전압과 전류가 주기적으로 커졌다 작아진다.
겉모습은 다르지만, 이 현상들은 모두 비슷한 구조를 가진다.

진동이 생기려면 보통 두 가지 성질이 필요하다.
첫째, 평형점으로 돌아가려는 복원 작용이 있어야 한다.
기계 시스템에서는 스프링이 이 역할을 하고, 전기 시스템에서는 커패시터에 쌓인 전하와 전압이 이 역할을 한다.
뒤에서 보겠지만 직렬 RLC 방정식에서 복원항은 $q/C$로 나타난다.
둘째, 바로 멈추지 않고 지나가게 만드는 관성 또는 에너지 저장 작용이 있어야 한다.
기계 시스템에서는 질량이 이 역할을 하고, 전기 시스템에서는 인덕터 전류와 자기장 에너지가 이 관성 같은 역할을 한다.
커패시터와 인덕터가 서로 에너지를 주고받기 때문에 전기적 진동이 만들어진다.

가장 단순한 진동 방정식은
\begin{equation}
  \ddot{x}+\omega_n^2 x=0
\end{equation}
이다.
여기서 $x$는 꼭 기계적 위치만 뜻하는 것이 아니라, 2차 시스템의 상태를 임시로 부르는 일반 기호이다.
기계 시스템에서는 위치가 될 수 있고, 전기 시스템에서는 전하나 전압처럼 시간에 따라 변하는 물리량이 될 수 있다.
이 식은 가속도가 위치와 반대 방향으로 생긴다는 뜻이다.
$x>0$이면 $\ddot{x}<0$이므로 왼쪽으로 가속되고, $x<0$이면 $\ddot{x}>0$이므로 오른쪽으로 가속된다.
그래서 시스템은 평형점 $x=0$ 주변을 왔다 갔다 한다.

이 방정식의 해는
\begin{equation}
  x(t)=A\cos(\omega_n t)+B\sin(\omega_n t)
\end{equation}
이다.
$\omega_n$은 얼마나 빠르게 흔들리는지를 정한다.
값이 크면 짧은 시간 안에 여러 번 흔들리고, 값이 작으면 천천히 흔들린다.
진동의 주기는
\begin{equation}
  T=\frac{2\pi}{\omega_n}
\end{equation}
이다.

현실의 진동은 대부분 시간이 지나면 줄어든다.
그 이유는 에너지가 조금씩 빠져나가기 때문이다.
기계 시스템에서는 마찰과 공기저항, 댐퍼가 에너지를 열로 바꾼다.
전기 시스템에서는 저항이 전기 에너지를 열로 바꾼다.
이처럼 에너지를 줄이는 작용을 감쇠(damping)라고 부른다.
감쇠가 없으면 진동은 계속되고, 감쇠가 있으면 진동은 점점 작아진다.
감쇠가 너무 크면 진동은 사라지지만 움직임도 느려진다.

따라서 진동을 이해할 때는 세 가지 질문을 계속 붙잡으면 된다.
\begin{itemize}
  \item 무엇이 평형점으로 되돌리려 하는가?
  \item 무엇이 바로 멈추지 않고 계속 지나가게 만드는가?
  \item 무엇이 에너지를 빼앗아 진동을 줄이는가?
\end{itemize}
이 세 질문에 대한 답이 각각 스프링, 질량, 댐퍼이고, 전기 회로에서는 커패시터, 인덕터, 저항으로 대응된다.

전기회로에서 이 대응이 처음에는 조금 낯설 수 있다.
기계 시스템에서는 물체가 왼쪽과 오른쪽으로 실제로 움직이는 모습이 눈에 보인다.
하지만 RLC 회로에서는 전하와 전류가 회로 안에서 오가고, 에너지가 전기장과 자기장 사이를 이동한다.
눈으로는 보이지 않지만 수식 구조는 매우 유사하다.
스프링이 변위에 따라 퍼텐셜 에너지를 저장하듯, 커패시터는 전압 또는 전하에 따라 전기장 에너지를 저장한다.
질량이 속도에 따라 운동 에너지를 저장하듯, 인덕터는 전류에 따라 자기장 에너지를 저장한다.
댐퍼가 운동 에너지를 열로 없애듯, 저항은 전기 에너지를 열로 없앤다.
이 장의 목표는 회로 해석을 완전히 끝내는 것이 아니라, 이 세 역할이 어떻게 2차 시스템의 진동을 만드는지 보는 것이다.

\subsection{저항이란 무엇인가?}

저항(resistor)은 전류가 지나가는 동안 전기 에너지의 일부를 열로 바꾸는 소자이다.
그래서 회로에서 가장 먼저 만나는 흐름의 어려움이 보통 저항이다.
전구의 필라멘트가 뜨거워지고, 전기히터가 열을 내고, 작은 저항 소자가 미세하게 데워지는 현상은 모두 같은 그림으로 볼 수 있다.

이 관계를 숫자로 잡아 준 사람이 게오르크 지몬 옴(Georg Simon Ohm)이다.
옴은 1820년대에 도선에 걸린 전압과 흐르는 전류 사이의 관계를 실험적으로 연구했고, 오늘날 우리가 옴의 법칙이라고 부르는 관계를 정리했다 \cite{ohm1827galvanischekette}.
\begin{equation}
  V=IR.
\end{equation}
여기서 $V$는 전압(volt, V), $I$는 전류(ampere, A), $R$은 저항(ohm, $\Omega$)이다.
이 식은 단순하지만 강력하다.
전압을 두 배로 올리면 전류도 두 배가 되고, 저항을 두 배로 키우면 같은 전압에서 전류는 절반으로 줄어든다.
즉 저항은 직류 회로에서 전류가 얼마나 쉽게 흐르는지, 또는 얼마나 어렵게 흐르는지를 한 숫자로 말해 준다.
예를 들어 $R=2~\Omega$인 저항에 $I=3~\mathrm{A}$가 흐르면 필요한 전압은 $V=6~\mathrm{V}$이고, 저항에서 소모되는 전력은 $p=VI=18~\mathrm{W}$이다.

물리적으로 보면 금속 안의 전자는 전기장에 의해 움직이지만, 격자 원자, 불순물, 열적 진동과 계속 상호작용한다.
전자는 전기장에서 에너지를 얻고, 그 에너지는 충돌과 산란을 통해 물질 내부의 열에너지로 바뀐다.
그래서 저항은 단순히 전자의 통로가 좁다는 뜻만은 아니다.
전자가 지나가면서 물질 내부와 계속 에너지를 주고받고, 그 결과 전기 에너지가 회수하기 어려운 열로 바뀐다는 뜻이다.
균일한 도선에서는 저항값이 대략
\begin{equation}
  R = \rho \frac{\ell}{A}
\end{equation}
로 주어진다.
여기서 $\rho$는 재료의 비저항(resistivity), $\ell$은 도선 길이, $A$는 단면적이다.
같은 재료라면 길수록 전자가 더 많은 산란 기회를 만나므로 저항이 커지고, 굵을수록 전류가 지나갈 길이 넓어져 저항이 작아진다.
구리선이 전선으로 많이 쓰이는 이유는 비저항이 작아 같은 길이와 굵기에서 전류를 비교적 잘 흘려보내기 때문이다.
수동부호약속으로 전류가 들어가는 쪽 단자를 전압의 $+$ 단자로 잡아 $v_R=Ri$라고 두면, 저항에서 소모되는 전력은
\begin{equation}
  p_R(t)=v_R(t)i(t)=Ri(t)^2
\end{equation}
이다.
여기서 소문자 $v_R(t)$, $i(t)$, $p_R(t)$는 시간에 따라 변하는 순간값이고, 앞의 대문자 $V$, $I$는 직류값이나 페이저처럼 한 상태를 대표하는 값으로 쓰인다.
전류의 방향이 바뀌어도 $i^2$는 항상 양수이므로, 저항은 에너지를 다시 돌려주지 않고 계속 열로 소산한다.
이 점에서 저항은 기계 시스템의 댐퍼와 닮았다.

이상적인 저항의 임피던스는
\begin{equation}
  Z_R=R
\end{equation}
이다.
여기서 임피던스는 전압과 전류의 비율을 주파수 영역에서 본 값이다.
허수부가 없다는 것은 에너지를 자기장이나 전기장에 저장하지 않는다는 뜻이다.
또 전압과 전류가 같은 위상이라는 뜻이기도 하다.
저항에는 인덕터 전류나 커패시터 전압처럼 오래 기억되는 저장 상태가 없다.
그래서 전류가 커지는 순간 전압도 같은 순간 커지고, 전류가 줄어드는 순간 전압도 같이 줄어든다.
이 기억 없음 때문에 이상적인 저항의 임피던스는 주파수와 상관없는 실수 $R$로 남는다.
저항의 본질은 반응이 얼마나 늦는가보다, 흐르는 동안 에너지를 얼마나 잃는가에 가깝다.

\subsection{인덕터란 무엇인가?}

인덕터(inductor)는 아주 단순하게 말하면 전선을 코일처럼 감아 놓은 소자이다.
전선에 전류가 흐르면 그 주변에는 자기장(magnetic field)이 생긴다.
전선을 한 번만 지나가게 두면 자기장이 약하게 생기지만, 전선을 여러 번 감아 코일로 만들면 각 감긴 부분이 만드는 자기장이 서로 더해진다.
그래서 인덕터는 전류가 흐를 때 자기장에 에너지를 저장하는 장치이다.
코일에 전류를 키운다는 것은 코일 주변 자기장을 키운다는 뜻이다.
자기장을 키우려면 에너지를 넣어야 하므로, 인덕터 전류는 공짜로 점프하지 않는다.

인덕터의 역사적 배경에는 전자기 유도(electromagnetic induction)가 있다.
마이클 패러데이(Michael Faraday)는 1831년에 변하는 자기선속이 유도기전력, 즉 전하를 밀어 움직이게 하는 전기적 효과를 만들 수 있음을 보였고, 이 연구를 1832년에 전자기 유도 실험으로 발표했다 \cite{faraday1832experimentalresearches}.
닫힌 도체 경로가 있으면 이 유도기전력 때문에 전류가 흐를 수 있다.
이 발견은 전기와 자기장이 서로 따로 노는 현상이 아니라 서로 깊게 연결되어 있음을 보여주었다.
이때 유도 효과는 원래 변화를 방해하는 부호로 나타난다.
이 부호 규칙을 흔히 렌츠의 법칙이라고 부른다.
인덕터가 전류 변화에 저항하는 이유가 바로 여기에 있다.
인덕터는 전류가 흐르는 것 자체를 막는 소자라기보다, 전류가 갑자기 바뀌려 할 때 더 큰 전압을 요구하는 소자라고 이해하면 좋다.
조금 더 정확히 말하면 전류 그 자체가 아니라 전류의 시간 변화에 대해 반대 방향의 유도기전력을 만든다.

이 성질은 다음 식으로 표현된다.
앞으로 소자 전압과 전류는 수동부호약속(passive sign convention)을 기준으로 쓴다.
수동부호약속이란 전류가 들어가는 쪽 단자를 전압의 $+$ 단자로 잡는 약속이다.
이 약속에서는 소자가 바깥 회로에서 에너지를 흡수할 때 전력 $p=vi$가 양수가 된다.
\begin{equation}
  v_L(t) = L\frac{\dd i(t)}{\dd t}.
  \label{eq:inductor_voltage}
\end{equation}
여기서 $v_L$은 인덕터 양단 전압, $i$는 인덕터에 흐르는 전류, $L$은 인덕턴스(inductance)이다.
단위는 헨리(H, henry)이다.
인덕턴스 $L$이 크다는 말은 같은 전류 변화율을 만들기 위해 더 큰 전압이 필요하다는 뜻이다.
코일을 더 많이 감거나, 자기장을 잘 모아 주는 철심 같은 코어를 넣으면 보통 인덕턴스가 커진다.
직관적으로는 전류 변화가 만들어야 하는 자기장 변화가 더 커져서, 회로가 그 변화를 더 무겁게 느낀다고 볼 수 있다.
식 \eqref{eq:inductor_voltage}를 말로 풀면 다음과 같다.

\begin{itemize}
  \item 전류가 천천히 바뀌면 $\dd i/\dd t$가 작으므로 필요한 전압도 작다.
  \item 전류가 빠르게 바뀌면 $\dd i/\dd t$가 크므로 큰 전압이 필요하다.
  \item 전류를 순간적으로 바꾸려면 $\dd i/\dd t$가 무한대가 되어야 하므로, 이상적으로는 무한대 전압이 필요하다.
\end{itemize}

그래서 이상적인 인덕터에서는 전류가 순간적으로 점프할 수 없다.
이것이 인덕터가 회로에서 만드는 첫 번째 방해이다.
전류가 빠르게 바뀌려는 순간, 인덕터 양단에는 그 변화를 만들어 내는 데 필요한 부호 있는 단자전압 $v_L=L\dd i/\dd t$가 생긴다.
전류가 증가할 때는 수동부호약속에서 양의 전압강하가 필요하고, 전류가 감소할 때는 반대 부호의 단자전압이 나타나며 인덕터가 저장한 에너지를 외부 회로에 돌려줄 수도 있다.
내부에서 보면 변하는 자기장이 원래 변화를 방해하는 방향의 유도기전력을 만든다고 말할 수 있다.
단자 전압강하와 반대 기준으로, 즉 전류 기준 방향으로 본 내부 유도기전력은 보통
\begin{equation}
  e_{\mathrm{ind}}(t)=-L\frac{\dd i(t)}{\dd t}
\end{equation}
처럼 부호를 붙여 쓴다.
즉 단자 전압강하 $v_L=L\dd i/\dd t$와 내부 유도기전력 $e_{\mathrm{ind}}=-L\dd i/\dd t$는 같은 물리 현상을 서로 반대 기준에서 본 표현이다.
부호가 부담스럽다면 첫 독해에서는 전류를 키울 때 에너지가 자기장으로 들어가고, 전류가 줄 때 저장된 에너지가 회로로 돌아올 수 있다는 점만 잡아도 충분하다.
인덕터의 방해는 단순한 마찰이 아니라 자기장에 에너지를 저장해야 하기 때문에 생기는 방해이다.
식 \eqref{eq:inductor_voltage}를 조금 바꾸어 쓰면 이 말이 더 직접적으로 보인다.
\begin{equation}
  \Delta i
  =
  \frac{1}{L}\int_{t_0}^{t_1} v_L(t)\,\dd t.
\end{equation}
전류를 얼마만큼 바꾸려면, 짧은 시간 동안 큰 전압을 주거나, 작은 전압을 오랫동안 주어야 한다.
즉 인덕터 전류는 전압의 순간값 하나로 바로 정해지는 것이 아니라, 전압을 시간 동안 얼마나 누적해서 걸었는지에 따라 바뀐다.
이 점이 인덕터가 전류 변화에 버티는 물리적 이유이다.
숫자로 보면 $L=1~\mathrm{H}$인 인덕터에서 전류를 $2~\mathrm{A/s}$의 기울기로 증가시키려면 $v_L=2~\mathrm{V}$가 필요하다.
반대로 전류가 이미 일정한 직류 정상상태라면 $\dd i/\dd t=0$이므로 이상적인 인덕터 전압은 $0~\mathrm{V}$이다.

인덕터에 저장된 에너지는 다음과 같다.
\begin{equation}
  E_L(t)=\frac{1}{2}L i(t)^2.
  \label{eq:inductor_energy}
\end{equation}
이 식이 어디서 나오는지도 직접 확인할 수 있다.
순간 전력은
\begin{equation}
  p_L(t)=v_L(t)i(t)
\end{equation}
이고, 식 \eqref{eq:inductor_voltage}를 대입하면
\begin{equation}
  p_L(t)=L i(t)\frac{\dd i(t)}{\dd t}.
\end{equation}
그런데
\begin{equation}
  \frac{\dd}{\dd t}\left(\frac{1}{2}L i(t)^2\right)
  =
  L i(t)\frac{\dd i(t)}{\dd t}
\end{equation}
이므로, 인덕터에 들어간 전력은 식 \eqref{eq:inductor_energy}의 에너지 변화율과 같다.
즉, 인덕터에 전력을 넣으면 열로 사라지는 것이 아니라 자기장 에너지로 저장된다.
전류가 줄어드는 구간에서는 $p_L$이 음수가 될 수 있다.
이때 인덕터는 저장해 두었던 자기장 에너지를 다시 회로로 돌려준다.
저항과 달리, 이상적인 인덕터는 에너지를 없애는 소자가 아니라 저장했다가 되돌려주는 소자이다.

이제 주파수 관점으로 보자.
정현파 전류 $i(t)=I\cos(\omega t)$를 인덕터에 흘리면 전류의 시간 미분은
\begin{equation}
  \frac{\dd i(t)}{\dd t}
  =
  -\omega I\sin(\omega t)
\end{equation}
이다.
주파수 $\omega$가 클수록 같은 전류 진폭을 만들기 위해 필요한 전압이 커진다.
또 $-\sin(\omega t)$는 $\cos(\omega t+90^\circ)$와 같으므로, 인덕터 전압은 전류보다 한 주기의 1/4만큼, 즉 $90^\circ$ 앞선다.
페이저 표기에서 시간 미분은 $\jj\omega$를 곱하는 일로 바뀌고, 여기서 $\jj$가 바로 $90^\circ$ 위상 선행을 나타낸다.
그래서 인덕터 임피던스는
\begin{equation}
  Z_L(\jj\omega)=\jj\omega L
\end{equation}
이고, 크기는
\begin{equation}
  |Z_L|=\omega L
\end{equation}
이다.
즉, 인덕터는 낮은 주파수에서는 거의 방해하지 않지만, 높은 주파수에서는 강하게 방해한다.
직류 정상상태에서는 전류가 더 이상 변하지 않으므로 $\dd i/\dd t=0$이고, 이상적인 인덕터는 마치 도선처럼 보인다.

\subsection{커패시터란 무엇인가?}

커패시터(capacitor)는 서로 떨어져 마주 보는 두 도체 사이에 전하를 분리해 두는 소자이다.
전체로 보면 중성이지만, 한쪽 판에는 $+q$, 다른 쪽 판에는 $-q$가 놓이고 그 사이에 전기장(electric field)이 생긴다.
커패시터가 저장하는 것은 전하 덩어리라기보다, 이렇게 분리된 전하가 만든 전기장 에너지라고 보는 편이 정확하다.
한국어 이름인 축전기는 그래서 꽤 좋은 이름이다.
전하를 쌓아 둔다는 직관을 주되, 실제로는 두 판 사이에 전하를 나누어 놓는다는 점까지 함께 떠올리게 해 주기 때문이다.

커패시터의 조상으로 자주 언급되는 장치는 라이덴 병(Leyden jar)이다.
1740년대 중반, 유리병 안팎에 도체를 두고 전하를 저장할 수 있다는 사실이 발견되었고, 이 장치는 초기 전기 실험에서 중요한 역할을 했다 \cite{leyden_jar}.
오늘날의 눈으로 보면 라이덴 병은 거대한 병 모양 커패시터이다.
유리가 절연체 역할을 하고, 안쪽과 바깥쪽 도체가 서로 다른 전하를 저장한다.
당시 사람들에게 라이덴 병은 전기가 순간적으로 지나가는 신기한 현상에 그치지 않고, 실제로 어딘가에 모아둘 수 있는 양이라는 사실을 보여준 장치였다.
커패시터라는 말이 직관적으로 어려울 수 있지만, 축전기라는 한국어 이름은 꽤 정확하다.
전하와 전기장 에너지를 축적하는 장치이기 때문이다.

커패시터의 기본 관계는
\begin{equation}
  q(t)=C v_C(t)
  \label{eq:capacitor_charge}
\end{equation}
이다.
여기서 $q$는 선택한 $+$ 판에 저장된 부호 있는 전하, $v_C$는 커패시터 양단 전압, $C$는 커패시턴스(capacitance)이다.
반대쪽 판에는 $-q$가 있으므로 이상적인 커패시터 전체는 여전히 전기적으로 중성이다.
단위는 패럿(F, farad)이다.
커패시턴스가 크다는 것은 같은 전압에서 더 많은 전하 분리를 견딜 수 있다는 뜻이다.
평행판 커패시터를 아주 단순하게 보면 커패시턴스는 대략
\begin{equation}
  C = \epsilon \frac{A}{d}
\end{equation}
로 주어진다.
여기서 $A$는 마주 보는 판의 면적, $d$는 판 사이 거리, $\epsilon$은 사이에 들어간 물질의 유전율(permittivity)이다.
판이 넓을수록 전하가 모일 자리가 많고, 판 사이가 가까울수록 같은 전하가 더 강하게 서로 끌어당긴다.
유전율이 큰 절연체를 넣어도 같은 전압에서 더 많은 전하 분리를 견딜 수 있다.
그래서 큰 커패시터는 전압 변화를 더 부드럽게 받아주는 전기적 저장소처럼 행동한다.

전류는 전하가 흐르는 속도이므로
\begin{equation}
  i_C(t)=\frac{\dd q(t)}{\dd t}
\end{equation}
이다.
여기서도 수동부호약속을 사용해, 전류가 커패시터의 $+$ 단자로 들어가면 $q$가 증가하는 방향으로 잡는다.
식 \eqref{eq:capacitor_charge}를 미분하면
\begin{equation}
  i_C(t)=C\frac{\dd v_C(t)}{\dd t}.
  \label{eq:capacitor_current}
\end{equation}
이 식을 말로 풀면 다음과 같다.

\begin{itemize}
  \item 커패시터 전압을 천천히 바꾸려면 작은 전류만 있으면 된다.
  \item 커패시터 전압을 빠르게 바꾸려면 큰 전류가 필요하다.
  \item 전압을 순간적으로 바꾸려면 $\dd v_C/\dd t$가 무한대가 되어야 하므로, 이상적으로는 무한대 전류가 필요하다.
\end{itemize}

그래서 이상적인 커패시터에서는 전압이 순간적으로 점프할 수 없다.
커패시터가 막는 것은 전류 그 자체라기보다 전압의 갑작스러운 변화이다.
전압을 빠르게 바꾸려면 두 판의 전하 분포와 전기장 에너지도 빠르게 바뀌어야 하고, 그러려면 그만큼 많은 전하가 짧은 시간에 이동해야 한다.
식 \eqref{eq:capacitor_current}도 누적 관점으로 다시 쓸 수 있다.
\begin{equation}
  \Delta v_C
  =
  \frac{1}{C}\int_{t_0}^{t_1} i_C(t)\,\dd t.
\end{equation}
커패시터 전압을 바꾸려면 전류를 흘려 전하를 실제로 옮겨야 한다.
큰 커패시턴스 $C$는 같은 전류를 같은 시간 동안 흘려도 전압이 조금만 바뀐다.
그래서 큰 커패시터는 전원 회로나 필터에서 흔들리는 전압을 완만하게 만드는 저장소처럼 쓰인다.
예를 들어 $C=1000~\mu\mathrm{F}=0.001~\mathrm{F}$인 커패시터에 $i_C=1~\mathrm{mA}=0.001~\mathrm{A}$를 계속 흘리면
\begin{equation}
  \frac{\dd v_C}{\dd t}=\frac{i_C}{C}=1~\mathrm{V/s}
\end{equation}
의 속도로 전압이 변한다.

커패시터에 저장된 에너지는
\begin{equation}
  E_C(t)=\frac{1}{2}C v_C(t)^2
  =
  \frac{q(t)^2}{2C}
  \label{eq:capacitor_energy}
\end{equation}
이다.
이 식도 전력에서 유도할 수 있다.
순간 전력은
\begin{equation}
  p_C(t)=v_C(t)i_C(t)
\end{equation}
이고, 식 \eqref{eq:capacitor_current}를 대입하면
\begin{equation}
  p_C(t)=C v_C(t)\frac{\dd v_C(t)}{\dd t}.
\end{equation}
이는
\begin{equation}
  \frac{\dd}{\dd t}\left(\frac{1}{2}C v_C(t)^2\right)
\end{equation}
와 같다.
따라서 커패시터에 들어간 전력은 전기장 에너지로 저장된다.
전압이 줄어드는 구간에서는 $p_C$가 음수가 될 수 있다.
이때 커패시터는 저장해 두었던 전기장 에너지를 회로로 되돌려준다.

주파수 관점에서 보면 커패시터의 성질은 인덕터와 반대이다.
정현파 전압 $v_C(t)=V\cos(\omega t)$를 커패시터에 걸면
\begin{equation}
  i_C(t)
  =
  C\frac{\dd v_C(t)}{\dd t}
  =
  -\omega C V\sin(\omega t)
  =
  \omega C V\cos(\omega t+90^\circ)
\end{equation}
가 된다.
즉 커패시터에서는 전류가 전압보다 $90^\circ$ 앞선다.
커패시터 임피던스는
\begin{equation}
  Z_C(\jj\omega)=\frac{1}{\jj\omega C}
\end{equation}
이고, 크기는
\begin{equation}
  |Z_C|=\frac{1}{\omega C}
\end{equation}
이다.
페이저로 한 줄만 더 쓰면, 식 $i_C=C\dd v_C/\dd t$는
\begin{equation}
  I_C=\jj\omega C V_C
\end{equation}
가 되므로
\begin{equation}
  Z_C=\frac{V_C}{I_C}
  =
  \frac{1}{\jj\omega C}
  =
  -\frac{\jj}{\omega C}
\end{equation}
가 된다.
전류가 전압보다 앞서기 때문에, 전압/전류 비율인 임피던스에는 음의 허수부가 나타난다.
주파수 $\omega$가 커질수록 커패시터 임피던스는 작아진다.
즉, 커패시터 가지(branch)에 전압을 걸고 그 가지 전류를 보면 빠르게 변하는 신호일수록 더 쉽게 전류가 흐르는 것처럼 보인다.
따라서 커패시터가 전압 변화에 버틴다는 말과 높은 주파수에서 전류가 잘 흐르는 것처럼 보인다는 말은 서로 반대가 아니다.
전압을 빠르게 흔들려면 많은 전하가 빠르게 오가야 하므로 큰 전류가 나타난다는 뜻이다.
다만 실제 회로에서 어떤 신호가 출력으로 통과하는지는 커패시터가 직렬로 놓였는지, 병렬로 접지 쪽에 놓였는지 같은 연결 방식에 따라 달라진다.
직류 정상상태에서는 전압이 일정하므로 $\dd v_C/\dd t=0$이고, 이상적인 커패시터에는 전류가 흐르지 않는다.
그래서 커패시터는 직류에서는 끊어진 회로처럼 보인다.

\subsection{인덕터와 커패시터는 왜 방해하는가?}

저항, 인덕터, 커패시터는 모두 회로의 흐름을 어렵게 만들지만, 어려움을 만드는 방식은 서로 다르다.
저항은 지나가는 전기 에너지를 열로 내보낸다.
한번 열로 흩어진 에너지는 이상적인 회로 해석에서는 다시 돌아오지 않는다.
반면 인덕터와 커패시터는 에너지를 없애기보다 잠시 맡아 둔다.
다만 에너지를 저장하려면 그 소자의 상태가 바뀌어야 하고, 유한한 전압, 전류, 전력으로는 저장 에너지를 순간적으로 무한히 바꿀 수 없다.
여기서 상태 변수(state variable)란 그 요소가 에너지를 저장하는 데 직접 쓰는 양이라고 생각하면 된다.
인덕터는 전류에, 커패시터는 전압 또는 전하에, 질량은 속도에, 스프링은 위치에 에너지를 저장한다.
여기서 저항한다는 말은 소자 이름인 저항 $R$과 같은 뜻이 아니다.
인덕터와 커패시터는 저항처럼 평균적으로 에너지를 열로 태워 없애는 것이 아니라, 저장 상태를 갑자기 바꾸기 어렵게 만들어 급격한 변화에 버티는 것처럼 보인다.

인덕터의 상태 변수는 전류이다.
인덕터 전류는 자기장 에너지와 직접 연결된다.
\begin{equation}
  E_L=\frac{1}{2}Li^2
\end{equation}
전류를 갑자기 바꾸려면 자기장 에너지도 갑자기 바뀌어야 한다.
그렇게 하려면 매우 큰 전압이 필요하다.
따라서 인덕터는 전류 변화에 저항한다.

커패시터의 상태 변수는 전압이다.
커패시터 전압은 전기장 에너지와 직접 연결된다.
\begin{equation}
  E_C=\frac{1}{2}Cv^2
\end{equation}
전압을 갑자기 바꾸려면 전기장 에너지도 갑자기 바뀌어야 하고, 두 판에 쌓인 전하도 갑자기 이동해야 한다.
그렇게 하려면 매우 큰 전류가 필요하다.
따라서 커패시터는 전압 변화에 저항한다.

이제 위상 특성도 자연스럽게 이해할 수 있다.
아래 설명은 정현파 정상상태에서, 전류가 수동 소자의 양의 단자로 들어간다고 보는 수동부호약속을 기준으로 한다.
인덕터에서는 전압이 먼저 걸려야 전류가 변하기 시작한다.
그래서 인덕터에서는 전압이 전류보다 $90^\circ$ 앞선다.
커패시터에서는 전류가 먼저 흘러 전하를 쌓아야 전압이 생긴다.
그래서 커패시터에서는 전류가 전압보다 $90^\circ$ 앞선다.
이 위상차가 복소수 임피던스에서 허수 단위 $\jj$로 표현된다.
비정현파 신호에서는 같은 말을 신호 전체에 한 번에 적용하기보다, 각 주파수 성분마다 위상 선행과 지연을 따로 본다.

정리하면 저항은 에너지를 잃게 만드는 소자이고, 인덕터와 커패시터는 저장된 에너지의 상태가 너무 빨리 바뀌지 못하게 하는 소자이다.
주파수 영역에서 이 서로 다른 반응을 한꺼번에 묶어 보는 언어가 임피던스이다.

역사적으로도 이 세 요소는 서로 다른 문제에서 등장했다.
옴은 전류와 전압의 비례 관계를 수학적으로 다룰 수 있게 만들었고, 패러데이는 변하는 자기장과 유도 전압의 관계를 실험적으로 드러냈으며, 라이덴 병은 전하를 분리해 저장할 수 있다는 사실을 보여준 초기 장치였다 \cite{ohm1827galvanischekette,faraday1832experimentalresearches,leyden_jar}.
따라서 RLC 회로를 외울 때도 단순히 세 알파벳을 나열하기보다, 표~\ref{tab:rlc-history-role}처럼 역사적 단서와 물리적 역할을 같이 붙여두면 훨씬 잘 기억된다.

\begin{table}[!htbp]
  \centering
  \caption{역사적 단서로 기억하는 RLC 소자의 역할}
  \small
  \setlength{\tabcolsep}{3pt}
  \begin{tabularx}{\linewidth}{@{}>{\raggedright\arraybackslash}p{0.13\linewidth}>{\raggedright\arraybackslash}p{0.25\linewidth}>{\raggedright\arraybackslash}p{0.23\linewidth}>{\raggedright\arraybackslash}X@{}}
    \toprule
    항목 & 역사적 단서 & 물리적 역할 & 임피던스 관점의 기억 문장 \\
    \midrule
    저항 $R$ & 옴의 회로 법칙: 전압과 전류의 비례 관계를 숫자로 다룬다. & 전기 에너지를 열로 소산한다. & 흐르는 동안 에너지를 잃게 만든다. 전기 쪽에서 댐퍼와 닮은 소산 요소이다. \\
    인덕터 $L$ & 패러데이의 전자기 유도: 변하는 자기장이 유도 전압을 만든다. & 전류와 연결된 자기장 에너지를 저장한다. & 전류가 갑자기 바뀌는 것을 어렵게 만든다. 전기 쪽에서 질량과 닮은 관성 요소이다. \\
    커패시터 $C$ & 라이덴 병: 떨어진 도체 사이에 전하를 분리해 저장한다. & 전압 또는 전하와 연결된 전기장 에너지를 저장한다. & 전압이 갑자기 바뀌는 것을 어렵게 만든다. 전기 쪽에서 스프링과 닮지만, 복원 계수에 해당하는 항은 $1/C$이다. \\
    RLC 회로 & 전신선, 교류 회로, 튜닝 문제에서 소산과 두 저장 요소를 함께 보게 된다. & $L$과 $C$가 에너지를 주고받고, $R$이 그 에너지를 조금씩 소산한다. & 저장소 두 개가 에너지를 왕복시키고, 저항이 매 주기마다 일부를 열로 빼앗는다. 그래서 진동, 공진, 감쇠가 한 식에 나타난다. \\
    \bottomrule
  \end{tabularx}
  \label{tab:rlc-history-role}
\end{table}

이를 상태 변화의 관점으로 다시 줄이면 표~\ref{tab:rlc-state-memory}처럼 볼 수 있다.
\begin{table}[!htbp]
  \centering
  \caption{RLC 소자를 상태 변화 관점에서 기억하기}
  \small
  \setlength{\tabcolsep}{3pt}
  \begin{tabularx}{\linewidth}{@{}>{\raggedright\arraybackslash}p{0.15\linewidth}>{\raggedright\arraybackslash}p{0.19\linewidth}>{\raggedright\arraybackslash}p{0.23\linewidth}>{\raggedright\arraybackslash}X@{}}
    \toprule
    소자 & 저장 또는 소산 & 순간적으로 바뀌기 어려운 양 & 기억할 말 \\
    \midrule
    저항 $R$ & 에너지 소산 & 저장 상태 없음 & 전류가 흐르는 동안 전기 에너지를 열로 바꾼다. 이상 모델에서는 전압과 전류가 즉시 비례한다. \\
    인덕터 $L$ & 자기장 에너지 저장 & 전류 $i$ & 전류를 바꾸려면 자기장 에너지를 바꾸어야 하므로 전압을 시간 동안 누적해서 걸어야 한다. \\
    커패시터 $C$ & 전기장 에너지 저장 & 전압 $v_C$ 또는 전하 $q$ & 전압을 바꾸려면 전하를 실제로 옮겨야 하므로 전류를 시간 동안 누적해서 흘려야 한다. \\
    \bottomrule
  \end{tabularx}
  \label{tab:rlc-state-memory}
\end{table}

이제 소자 하나에서 회로 전체로 넘어간다.
저항은 전류가 흐르는 동안 에너지를 소산하고, 인덕터는 전류 상태가 갑자기 바뀌는 데 맞서며, 커패시터는 전하와 전압 상태가 갑자기 바뀌는 데 맞선다.
직렬 RLC에서는 같은 전류가 세 소자를 지나가므로, 다음 방정식의 $L\ddot{q}$, $R\dot{q}$, $q/C$를 각각 전류 변화에 대한 전압, 전류 소산 전압, 전하가 만든 복원 전압으로 읽으면 된다.

\subsection{직렬 RLC 회로}

직렬 RLC 회로는 저항 $R$, 인덕터 $L$, 커패시터 $C$가 하나의 전류 경로에 놓인 회로이다.
회로를 한 방향으로 도는 기준 전류를 $i(t)$라고 잡고, 그 전류가 들어가는 커패시터 판의 전하를 $q(t)$라고 하자.
그러면 전류는
\begin{equation}
  i(t) = \dot{q}(t)
\end{equation}
이다.
키르히호프 전압 법칙(Kirchhoff voltage law)은 닫힌 회로 한 바퀴를 돌며 만나는 전압 변화의 대수합이 0이라는 법칙이다.
전압은 전하 $1~\mathrm{C}$당 에너지 장부처럼 볼 수 있다.
전원은 장부에 에너지를 더하고, 세 소자는 각각 열, 자기장, 전기장 형태로 그 에너지를 가져간다.
여기서는 전원 전압이 상승이고 각 소자 전압이 하강이 되도록 방향을 잡으면
\begin{equation}
  v_{\mathrm{in}}(t)-v_L(t)-v_R(t)-v_C(t)=0
\end{equation}
으로 쓸 수 있다.
직렬 RLC에서는 같은 전류가 저항, 인덕터, 커패시터를 차례로 지나가므로, 각 소자 전압을 더해 입력 전압과 맞춘다.
이를 적용하면
\begin{align}
  v_L(t) &= L\frac{\dd i(t)}{\dd t}=L\ddot{q}(t), \\
  v_R(t) &= Ri(t)=R\dot{q}(t), \\
  v_C(t) &= \frac{q(t)}{C}
\end{align}
이고, 직렬 회로에서는 세 소자 전압의 합이 입력 전압과 같다.
선택한 $+$ 판에 $q>0$의 전하가 쌓이면 커패시터 전압도 그 전하를 더 밀어 넣기 어렵게 만드는 방향으로 생긴다.
그래서 $q/C$는 단순히 전압 이름 하나가 아니라, 전하가 기준 상태에서 벗어난 만큼 커지는 전기적 effort 항처럼 읽을 수 있다.
따라서
\begin{equation}
  L \ddot{q}(t) + R \dot{q}(t) + \frac{1}{C} q(t) = v_{\mathrm{in}}(t)
\end{equation}
을 얻는다.
이 한 줄에는 세 소자의 방해 방식이 모두 들어 있다.
\begin{itemize}
  \item $L\ddot{q}$는 전류 $i=\dot{q}$를 바꾸려는 것에 대한 방해이다. 전류 변화가 빠를수록 큰 인덕터 전압이 필요하다.
  \item $R\dot{q}$는 전류가 흐르는 동안 생기는 소산이다. 전류가 클수록 저항에서 더 많은 에너지가 열로 빠져나간다.
  \item $\frac{1}{C}q$는 쌓인 전하와 커패시터 전압이 기준 상태에서 벗어난 만큼 커지는 전기적 effort 항이다. 이 항이 기계 시스템의 스프링 항 $kx$와 같은 자리에 놓인다.
\end{itemize}
여기서 force-voltage 비유란 전압을 힘에, 전류를 속도에 대응시키는 읽기 방식이다.
이 비유에서는 $L\leftrightarrow m$, $R\leftrightarrow b$, $\frac{1}{C}\leftrightarrow k$로 읽는 편이 가장 안전하다.
이는 2차 미분방정식이다.
이제 이 회로가 입력 전압을 받았을 때 전하나 커패시터 전압이 시간에 따라 어떻게 반응하는지 보고 싶다.
라플라스 변환은 같은 회로를 전달함수와 주파수 언어로 다시 쓰기 위한 계산 도구이다.
물리적 핵심은 바로 뒤의 에너지 해석에서 다시 확인한다.
이를 위해 먼저 라플라스 변환을 사용한다.
영 초기조건은 시작 순간에 인덕터 전류와 커패시터 전압 같은 저장 상태가 0이라고 두는 계산 조건이다.
영 초기조건에서 라플라스 변환을 적용하면
\begin{equation}
  \left(Ls^2 + Rs + \frac{1}{C}\right) Q(s) = V_{\mathrm{in}}(s)
\end{equation}
이다.
여기서 $Q(s)$는 전하 $q(t)$의 라플라스 변환이고, $V_{\mathrm{in}}(s)$는 입력 전압의 라플라스 변환이다.
전하에 대한 전달함수는
\begin{equation}
  \frac{Q(s)}{V_{\mathrm{in}}(s)}
  =
  \frac{1}{Ls^2 + Rs + \frac{1}{C}}
\end{equation}
이다.
전하 $q$를 먼저 출력으로 잡은 이유는 $i=\dot{q}$이고 $v_C=q/C$라서 전류와 커패시터 전압을 모두 쉽게 얻을 수 있기 때문이다.
예를 들어 커패시터 전압을 출력으로 잡으면
\begin{equation}
  \frac{V_C(s)}{V_{\mathrm{in}}(s)}
  =
  \frac{1}{LCs^2+RCs+1}
\end{equation}
이 된다.
전류를 출력으로 잡으면 $I(s)=sQ(s)$이므로
\begin{equation}
  \frac{I(s)}{V_{\mathrm{in}}(s)}
  =
  \frac{s}{Ls^2+Rs+\frac{1}{C}}
\end{equation}
로 쓸 수 있다.
즉 같은 RLC 회로라도 무엇을 출력으로 보느냐에 따라 전달함수의 모양이 달라진다.

이 방정식은 에너지 관점에서도 해석할 수 있다.
회로에 저장된 총 에너지는 인덕터의 자기장 에너지와 커패시터의 전기장 에너지의 합이다.
\begin{equation}
  E(t)
  =
  \frac{1}{2}L i(t)^2
  +
  \frac{1}{2C}q(t)^2.
\end{equation}
양변을 시간에 대해 미분하고 $i=\dot{q}$를 사용하면
\begin{equation}
  \dot{E}(t)
  =
  i(t)
  \left(
  L\dot{i}(t)+\frac{1}{C}q(t)
  \right).
\end{equation}
키르히호프 전압 법칙에서
\begin{equation}
  L\dot{i}(t)+\frac{1}{C}q(t)
  =
  v_{\mathrm{in}}(t)-Ri(t)
\end{equation}
이므로
\begin{equation}
  \dot{E}(t)
  =
  v_{\mathrm{in}}(t)i(t) - R i(t)^2.
\end{equation}
첫 번째 항은 외부 전원이 회로에 넣어주는 전력이고, 두 번째 항은 저항에서 열로 소산되는 전력이다.
인덕터와 커패시터는 서로 에너지를 주고받지만, 이상적인 경우 스스로 에너지를 없애지는 않는다.
따라서 RLC 회로의 진동은 인덕터와 커패시터 사이의 에너지 교환이며, 저항은 그 진동을 감쇠시키는 역할을 한다.
외부 구동을 0으로 두면 $v_{\mathrm{in}}=0$이므로 $\dot{E}=-Ri^2$이다.
저항도 없으면 $\dot{E}=0$이고, 에너지는 사라지지 않은 채 $L$과 $C$ 사이에서만 오간다.

\subsection{표준 2차 시스템과의 비교}

이 절의 $x$는 기계 위치가 아니라 두 번 미분되는 대표 변수 자리를 뜻한다.
RLC 회로에서는 그 자리에 전하 $q$가 들어가고, 다음 장의 기계 시스템에서는 실제 위치 $x$가 들어간다.

표준 2차 시스템은 보통 다음과 같이 쓴다.
\begin{equation}
  \ddot{x} + 2\zeta \omega_n \dot{x} + \omega_n^2 x = \omega_n^2 u.
\end{equation}
여기서 $\omega_n$은 고유진동수(natural frequency), $\zeta$는 감쇠비(damping ratio)이다.
이 표준형은 단순한 수학 정리가 아니라, 서로 다른 물리 시스템을 같은 눈금으로 비교하기 위한 약속이다.
질량-스프링-댐퍼나 RLC 회로처럼 1차원으로 쓸 수 있는 시스템에서는 방정식을 이 형태로 바꾸면 두 숫자가 바로 보인다.
나중에 볼 로봇 끝단의 가상 임피던스도 각 축을 독립적인 2차 시스템처럼 단순화해서 볼 때는 같은 방식으로 읽을 수 있다.
다만 이것은 뒤에서 다시 볼 미리보기이고, 실제 3차원 또는 6차원 Cartesian impedance에서는 $\bm{M}_d$, $\bm{D}_d$, $\bm{K}_d$ 행렬의 결합과 방향별 모드까지 함께 보아야 한다.
$\omega_n$은 시스템이 가진 자연스러운 흔들림 속도이고, $\zeta$는 그 흔들림이 얼마나 빨리 사라지는지를 나타내는 비율이다.
조금 더 정확히 말하면 $\zeta$는 임계감쇠와 비교한 상대 감쇠이고, 초 단위로 얼마나 빨리 줄어드는지는 $\zeta\omega_n$도 함께 보아야 한다.
같은 $\zeta$라도 $\omega_n$이 큰 회로는 포락선이 더 빠르게 줄 수 있고, 같은 RLC 직렬 회로에서는 $\zeta\omega_n=R/(2L)$로 나타난다.
그래서 복잡한 회로의 세부 부품값을 모두 보지 않아도, $\omega_n$과 $\zeta$만 보면 응답의 큰 윤곽을 먼저 예상할 수 있다.
지금은 두 이름만 들고 가도 된다.
다음 장에서 같은 값이 실제로 흔들리는 질량-스프링-댐퍼 그래프로 다시 등장한다.
RLC 회로 방정식
\begin{equation}
  L\ddot{q}+R\dot{q}+\frac{1}{C}q=v_{\mathrm{in}}
\end{equation}
을 표준형과 비교하려면 먼저 양변을 $L$로 나눈다.
\begin{equation}
  \ddot{q}+\frac{R}{L}\dot{q}+\frac{1}{LC}q=\frac{1}{L}v_{\mathrm{in}}.
\end{equation}
이제 표준형의 $x$를 전하 $q$로, 입력 $u$를 전압 입력에 비례한 값으로 생각하면 계수 비교가 가능해진다.
왼쪽의 각 항을 표준형과 나란히 놓으면 다음과 같다.
\begin{align}
  \text{속도항:}\quad 2\zeta\omega_n
  &=
  \frac{R}{L}, \\
  \text{위치항:}\quad \omega_n^2
  &=
  \frac{1}{LC}, \\
  \text{입력항:}\quad \omega_n^2 u
  &=
  \frac{1}{L}v_{\mathrm{in}}.
\end{align}
마지막 줄에서 $\omega_n^2=1/(LC)$이므로 표준형 입력은
\begin{equation}
  u=Cv_{\mathrm{in}}
\end{equation}
처럼 전압에 커패시턴스를 곱한 전하 단위의 입력으로 해석할 수 있다.
즉 표준형의 $u$는 실제 전원 장치가 내는 전압을 그대로 복사한 값이 아니라, 서로 다른 2차 시스템을 같은 모양으로 비교하려고 스케일을 바꾼 계산용 입력이다.
따라서 이 표준형에서 $u$의 단위는 전압이 아니라 전하 단위이다.
표준형으로 바꾸는 일은 물리 시스템을 바꾸는 것이 아니라 눈금을 바꾸는 일이다.
전하, 위치, 힘처럼 단위가 다른 시스템도 $\omega_n$과 $\zeta$라는 두 눈금으로 옮겨 놓으면 얼마나 빠르게 흔들리는가와 얼마나 빨리 사라지는가를 같은 그래프 언어로 비교할 수 있다.
$u=Cv_{\mathrm{in}}$ 같은 입력 스케일링도 실제 전압을 다른 장치로 바꾼다는 뜻이 아니라, 비교하기 쉬운 기준 입력으로 다시 표현하는 과정이다.
RLC 회로를 이 표준형과 비교하면
\begin{align}
  \omega_n &= \frac{1}{\sqrt{LC}}, \\
  \zeta &= \frac{R}{2} \sqrt{\frac{C}{L}}.
\end{align}
따라서 저항 $R$을 증가시키면 감쇠비가 커지고, 인덕턴스 $L$과 커패시턴스 $C$는 고유진동수와 에너지 저장 특성을 바꾼다.
에너지 관점에서 보면 $\omega_n$은 인덕터의 자기장 에너지와 커패시터의 전기장 에너지가 서로 주고받는 리듬이고, $\zeta$는 그 저장 리듬에 대해 저항 손실이 얼마나 큰지를 정규화한 비율이다.
전류가 흐를 때 저항은 $Ri^2$의 일률로 에너지를 열로 빼앗고, $L$과 $C$는 그 사이에서 에너지를 저장했다가 되돌려준다.
그래서 같은 저항값이라도 $L$과 $C$가 바뀌면 에너지 저장량과 교환 속도가 달라져 감쇠비의 의미도 함께 달라진다.
숫자로 한 번 계산해 보자.
$L=10~\mathrm{mH}=0.01~\mathrm{H}$, $C=100~\mu\mathrm{F}=0.0001~\mathrm{F}$이면
\begin{equation}
  \omega_n=\frac{1}{\sqrt{LC}}=1000~\mathrm{rad/s},
  \qquad
  f_n=\frac{\omega_n}{2\pi}\approx159~\mathrm{Hz}
\end{equation}
이다.
이 회로에서 임계감쇠에 해당하는 저항은 $\zeta=1$을 넣어
\begin{equation}
  R_c = 2\sqrt{\frac{L}{C}} = 20~\Omega
\end{equation}
로 계산된다.
$R$이 이보다 작으면 부족감쇠 응답이 나타나고, 이보다 크면 과감쇠 응답 쪽으로 간다.
전기 회로의 과도응답은 이 세 파라미터가 만드는 2차 시스템의 응답으로 이해할 수 있다.
다음 장의 식과 미리 나란히 놓으면 대응이 더 선명하다.
\begin{equation}
  L\ddot{q}+R\dot{q}+\frac{1}{C}q=v_{\mathrm{in}},
  \qquad
  m\ddot{x}+b\dot{x}+kx=f.
\end{equation}
여기서 전하 $q$는 기계 위치 $x$에, 전류 $i=\dot{q}$는 속도 $\dot{x}$에 대응된다.
관성 역할은 $L \leftrightarrow m$, 소산 역할은 $R \leftrightarrow b$, 기준 상태에서 벗어난 양에 대응하는 effort 역할은 $\frac{1}{C} \leftrightarrow k$로 대응된다.
다시 말해 커패시터 $C$ 자체가 스프링 강성 $k$와 같은 숫자라는 뜻이 아니라, 커패시터의 역수 $1/C$가 스프링 강성처럼 전하 변위에 비례하는 계수로 나타난다는 뜻이다.

전달함수에서 보았던 분모 \(Ls^2+Rs+1/C\)는 단순한 계산 결과가 아니다.
전달함수 절에서 본 \(Q(s)/V_{\mathrm{in}}(s)\)와 이제 볼 자유진동은 출발 조건은 다르지만, 같은 회로 방정식의 같은 왼쪽 항을 본다.
zero-state 전달함수는 저장 에너지를 0으로 두고 입력 전압이 어떤 전하 응답을 만드는지 묻는 실험이고, zero-input 자유응답은 입력 전압을 0으로 두고 처음 저장된 에너지가 어떤 모드로 사라지거나 흔들리는지 묻는 실험이다.
zero-state와 zero-input은 다른 실험이지만 같은 분모를 공유한다.
따라서 \(Ls^2+Rs+1/C\)의 근이 전달함수의 pole이면서 동시에 자유진동의 감쇠와 흔들림을 정한다.
여기서 pole은 처음 흔들린 뒤 어떤 속도로 사라지거나 계속 흔들리는지를 정하는 분모의 뿌리 정도로 읽으면 된다.
입력과 출력 선택은 분자와 관찰 신호를 바꾸지만, \(L\), \(R\), \(C\)가 만드는 자연 모드는 그대로 다음 장의 \(m\), \(b\), \(k\) 해석으로 옮겨 간다.

\subsection{전기적 진동과 공진}

여기부터는 외부 입력 때문에 생기는 zero-state 응답이 아니라, 처음 저장된 에너지만으로 움직이는 zero-input 자유응답을 먼저 본다.
외부 구동을 0으로 두되 회로 루프는 닫힌 상태로 두고, 충전된 커패시터와 인덕터 전류가 서로 에너지를 주고받는 장면을 보는 것이다.

저항이 없는 이상적인 LC 회로를 생각하면 방정식은
\begin{equation}
  L\ddot{q}(t)+\frac{1}{C}q(t)=0
\end{equation}
이 된다.
이는 단순 조화진동자(simple harmonic oscillator)와 같은 형태이다.
해는
\begin{equation}
  q(t)=A\cos(\omega_n t)+B\sin(\omega_n t)
\end{equation}
이고, 고유진동수는
\begin{equation}
  \omega_n=\frac{1}{\sqrt{LC}}
\end{equation}
이다.
이때 에너지는 커패시터의 전기장 에너지와 인덕터의 자기장 에너지 사이를 주기적으로 오간다.
커패시터 전압이 최대일 때는 전기장 에너지가 크고 전류는 작다.
전류가 최대일 때는 인덕터의 자기장 에너지가 크고 커패시터 전압은 작다.

이 장면을 아주 천천히 멈춰서 보면 다음 장의 기계 진동과 거의 같은 그림이 보인다.
처음에 커패시터가 충분히 충전되어 전압이 최대라고 하자.
이 순간에는 아직 전류가 거의 흐르지 않으므로 에너지가 주로 커패시터의 전기장에 저장되어 있다.
카트-스프링 시스템으로 바꾸어 말하면, 카트가 오른쪽 끝까지 밀려 있고 속도는 0인 순간과 같다.
그때 에너지는 질량의 운동 에너지가 아니라 스프링의 탄성 퍼텐셜 에너지에 저장되어 있다.
시간이 조금 지나면 커패시터가 방전되면서 전류가 커지고, 에너지가 인덕터의 자기장으로 이동한다.
기계 시스템에서는 카트가 평형점을 지나며 속도가 최대가 되고, 에너지가 질량의 운동 에너지로 이동하는 장면에 해당한다.
다시 시간이 지나면 전류가 줄어들고 반대 극성으로 커패시터가 충전된다.
카트도 반대쪽 끝까지 이동한 뒤 잠시 멈춘다.
즉, 전기적 진동은 전기장과 자기장 사이의 에너지 왕복이고, 기계적 진동은 스프링 퍼텐셜 에너지와 질량 운동 에너지 사이의 에너지 왕복이다.

저항 $R$이 추가되면 매 주기마다 일부 에너지가 열로 사라지므로 진폭이 줄어든다.
부족감쇠 조건 $0<\zeta<1$에서 감쇠 진동수는
\begin{equation}
  \omega_d = \omega_n\sqrt{1-\zeta^2}
\end{equation}
이다.
여기서 두 주파수를 나누어 읽자.
$\omega_d$는 외부 구동을 0으로 둔 닫힌 회로가 처음 저장된 에너지만으로 흔들릴 때, 전압이나 전류 플롯의 피크 간격에서 보이는 자유진동 속도이다.
반면 이어서 나오는 $\omega$는 밖에서 정현파 전압을 계속 넣을 때 우리가 조절하는 입력 주파수이다.
저항이 있으면 자유진동의 봉우리는 $\omega_d$에 가까운 간격으로 줄어들지만, 강제진동의 공진은 입력 주파수 $\omega$를 바꾸어 보며 응답 크기가 어디서 커지는지 찾는 문제이다.
따라서 고유진동수 근처라는 말은 입력 타이밍이 회로 내부의 에너지 교환 리듬과 가까워진다는 뜻이지, 모든 플롯에서 같은 주파수 기호 하나를 그대로 읽으라는 뜻은 아니다.
저항이 작을 때 외부 입력 주파수가 회로의 고유진동수 근처에 오면, 직렬 회로의 전체 임피던스는
\begin{equation}
  Z_{\mathrm{series}}(\jj\omega)
  =
  R+\jj\omega L+\frac{1}{\jj\omega C}
  =
  R+\jj\left(\omega L-\frac{1}{\omega C}\right)
\end{equation}
로 읽을 수 있다.
여기서 괄호 안의 $\omega L-1/(\omega C)$는 저항처럼 에너지를 태워 없애는 항이 아니라, 전류와 $90^\circ$ 어긋나 나타나는 순리액턴스이다.
인덕터는 $X_L=+\omega L$, 커패시터는 $X_C=-1/(\omega C)$를 갖는다.
이 부호는 전압과 전류의 위상 방향을 나타낸다.
따라서
\begin{equation}
  \omega L-\frac{1}{\omega C}=0
\end{equation}
이 되는 순간에는 전원에서 본 순허수 방해가 0이 되고, 전체 임피던스가 거의 $R$만 남는다.
이것이 상쇄의 뜻이다.
인덕터와 커패시터 양단 전압이 각각 사라진다는 뜻이 아니라, 두 전압이 서로 반대 위상으로 크게 나타나지만 전원 입장에서는 그 합이 작게 보인다는 뜻이다.
그래서 저항이 작으면 같은 입력 전압에서도 큰 전류가 흐를 수 있다.
이를 직렬 전류 공진(series current resonance)이라고 부를 수 있다.
전류 전달함수로 쓰면
\begin{equation}
  \frac{I(s)}{V_{\mathrm{in}}(s)}
  =
  \frac{s}{Ls^2+Rs+\frac{1}{C}}
\end{equation}
이고, 커패시터 전압을 출력으로 보면
\begin{equation}
  \frac{V_C(s)}{V_{\mathrm{in}}(s)}
  =
  \frac{1}{LCs^2+RCs+1}
\end{equation}
이다.
따라서 공진 근처에서 무엇이 크게 보이는지는 전류를 보는지, 커패시터 전압을 보는지, 인덕터 전압을 보는지에 따라 조금 다르게 표현된다.
이 글에서는 먼저 직렬 회로에서 전압 입력에 대한 회로 전류 응답이 커지는 현상을 중심으로 이해하면 충분하다.
이 현상을 공진이라고 부른다.
공진은 단순히 크게 흔들린다는 뜻이 아니다.
외부에서 에너지를 넣는 타이밍이 시스템 내부의 에너지 교환 리듬과 맞아떨어진다는 뜻이다.
그네를 밀 때도 아무 때나 미는 것보다 돌아오는 리듬에 맞춰 밀면 작은 힘으로도 크게 흔들린다.
RLC 회로에서도 입력 주파수가 LC 에너지 교환의 리듬과 가까우면 과도 구간에서 저장 에너지가 축적되고, 이후 정현파 정상상태에서는 평균 입력 전력이 저항 손실과 균형을 이루는 큰 전류 또는 전압 진폭이 형성될 수 있다.
따라서 공진은 위험한 현상이기도 하지만, 라디오 튜닝이나 필터처럼 원하는 주파수를 선택하는 데 활용되는 현상이기도 하다.

이 장의 역사적 조각들을 한 줄로 이어 보면 흐름이 보인다.
라이덴 병은 전하와 전기장 에너지를 저장할 수 있음을 보여주었고, 옴의 법칙은 저항을 숫자로 다룰 수 있게 했으며, 패러데이의 전자기 유도는 전류 변화와 자기장 변화가 서로 얽혀 있음을 보여주었다.
그리고 교류와 통신선 문제를 다루려면 이 세 요소를 임피던스와 주파수 응답이라는 한 언어로 묶어야 했다.

다음 장에서는 같은 2차 구조를 눈에 보이는 운동으로 다시 읽는다.
목표는 전기 부품을 기계 부품으로 억지 번역하는 것이 아니라, 같은 2차 미분방정식 구조가 눈에 보이는 운동에서는 어떻게 나타나는지 확인하는 것이다.
여기서 표준형의 $x$는 실제 기계 위치라는 뜻이 아니라, 2차 시스템에서 두 번 미분되는 대표 상태를 임시로 부르는 이름이다.
RLC 회로에서는 그 자리에 전하 $q$가 들어가고, 기계 시스템에서는 위치 $x$가 들어간다.
따라서 전류 $i=\dot{q}$는 위치 $x$가 아니라 속도 $\dot{x}$에 대응된다.
인덕터와 커패시터가 자기장과 전기장 사이에서 에너지를 주고받았다면, 질량-스프링 시스템은 운동 에너지와 탄성 퍼텐셜 에너지 사이에서 에너지를 주고받는다.
저항이 전기적 진동을 줄였듯이, 댐퍼는 기계적 진동을 줄인다.

이 장의 학습 산출물은 세 가지이다.
첫째, $R$, $L$, $C$가 각각 무엇을 소산하거나 저장하는지 말로 구분한다.
둘째, $L$, $R$, $C$가 주어졌을 때 $\omega_n=1/\sqrt{LC}$와 $\zeta=(R/2)\sqrt{C/L}$를 계산하고, $R$을 키우면 감쇠가 커지며 $L$과 $C$를 바꾸면 에너지 교환 리듬이 바뀐다고 예측한다.
셋째, 직렬 RLC와 질량-스프링-댐퍼를 나란히 볼 때 $L\leftrightarrow m$, $R\leftrightarrow b$, $1/C\leftrightarrow k$로 읽는다.
커패시턴스 $C$ 자체가 스프링 강성이라는 뜻이 아니라, $q/C$가 $kx$처럼 기준 상태에서 벗어난 양에 대응하는 effort 항 자리에 놓인다는 뜻이다.
